% About — King James Version
\clearpage
\thispagestyle{empty}%
\null\vspace*{16mm}
\begin{center}
  {\Large\booktitlefont\scshape About This Edition\par}
  \vspace{3pt}
  {\color{gospels!70}\rule{28mm}{0.4pt}\enspace{\char"2766}\enspace\rule{28mm}{0.4pt}\par}
\end{center}
\vspace{9mm}
\begin{center}\begin{minipage}{0.74\textwidth}
\setlength{\parskip}{6pt plus 2pt}\small

{\booktitlefont\scshape The King James Version}\par
The King James, or Authorized, Version was commissioned by King James~VI \& I at
the Hampton Court Conference of 1604 and completed in 1611. Some forty-seven
translators, working in six companies at Westminster, Oxford, and Cambridge,
rendered the Hebrew Masoretic Text and the Greek \emph{Textus Receptus} into
English, drawing deeply on the earlier work of William Tyndale and the Bishops'
Bible. Following early practice the text carries no quotation marks; the spelling
here follows the standardized 1769 Oxford text.

\medskip
{\booktitlefont\scshape Copyright}\par
The King James Version (Authorized Version, 1611) is in the public domain in the
United States and most of the world. In the United Kingdom the rights are vested
in the Crown and administered by Cambridge University Press, the King's Printer;
reproduction there is permitted only by royal letters patent.

% Shared colophon: how the edition was produced + the project's goals.
% Input by each edition's about.tex, inside a centred minipage.

\medskip
{\booktitlefont\scshape How this edition was made}\par
This volume is typeset with LuaLaTeX and David Purton's \texttt{scripture}
package, set in Georg Duffner's EB~Garamond on a strict baseline grid. The
scripture text is drawn from its publisher's source, converted to
scripture-package markup by the project's Python generators, and composed in the
manner of the early printed Bibles: hanging (margin-protruding) verse numbers,
drop-cap lettrine chapter openings, the divine name in small capitals, and
period public-domain woodcut headpieces vectorized from early books and assigned
by canonical division. Proper names are hyphenated after a published
Bible-typesetting authority, and widows and orphans are resolved without
disturbing the baseline grid.

\medskip
{\booktitlefont\scshape About the project}\par
\emph{geneve\_1564} is an open typesetting project in the visual tradition of the
Geneva Bible of 1564. Its goal is to render Scripture as a beautiful, readable
book --- honouring historic printing practice --- that anyone can build for
themselves, for print or for e-ink readers. It began as a fork of Rapha\"el
Pinson's Geneva~1564 reproduction and was inspired by the \texttt{scripture}
package. Source, provenance, and build instructions:
\texttt{github.com/coopermj/geneve\_1564}.

\end{minipage}\end{center}
\vfill
\clearpage
